% Figure: Energy Ratchet — Discrete Event Model
% Chapter 31 — Temporal Evolution (Section: The Energy Ratchet)
% Shows B(t) trajectory with ceiling B_max declining at each event,
% partial recovery between events, and severity thresholds.

\begin{figure}[htbp]
\centering
\begin{tikzpicture}[
    scale=0.82, every node/.style={scale=0.82},
    >=Stealth,
    every node/.append style={font=\sffamily},
]

%% ====================================================================
%% (a) Ratchet Mechanism — Single Patient Trajectory
%% ====================================================================
\begin{scope}[shift={(0,0)}]

  % --- Axes ---
  \draw[->, thick] (-0.3, 0) -- (11.5, 0)
    node[below, font=\sffamily\small] {Time};
  \draw[->, thick] (0, -0.3) -- (0, 7.8)
    node[above, rotate=90, anchor=south, font=\sffamily\small]
    {Functional Capacity $B(t)$};

  % --- Y-axis labels ---
  \node[left, font=\sffamily\scriptsize] at (-0.15, 7.0) {1.0};
  \node[left, font=\sffamily\scriptsize] at (-0.15, 0.3) {0};

  % --- Severity threshold bands ---
  % Mild zone (top)
  \fill[green!8] (0, 4.2) rectangle (11.2, 7.2);
  % Moderate zone
  \fill[yellow!8] (0, 2.45) rectangle (11.2, 4.2);
  % Severe zone
  \fill[orange!6] (0, 1.05) rectangle (11.2, 2.45);
  % Very severe zone (bottom)
  \fill[red!6] (0, 0) rectangle (11.2, 1.05);

  % --- Threshold lines ---
  \draw[green!50!black, densely dashed, thin] (0, 4.2) -- (11.2, 4.2);
  \draw[yellow!60!black, densely dashed, thin] (0, 2.45) -- (11.2, 2.45);
  \draw[red!50!black, densely dashed, thin] (0, 1.05) -- (11.2, 1.05);

  % --- Threshold labels (right side) ---
  \node[font=\sffamily\tiny, green!50!black, anchor=west] at (11.4, 4.2)
    {$\theta_{\mathrm{mod}}$};
  \node[font=\sffamily\tiny, yellow!60!black, anchor=west] at (11.4, 2.45)
    {$\theta_{\mathrm{sev}}$};
  \node[font=\sffamily\tiny, red!50!black, anchor=west] at (11.4, 1.05)
    {$\theta_{\mathrm{vs}}$};

  % --- Severity zone labels (left margin) ---
  \node[font=\sffamily\tiny, green!40!black, rotate=90] at (-0.7, 5.7) {Mild};
  \node[font=\sffamily\tiny, yellow!50!black, rotate=90] at (-0.7, 3.3) {Moderate};
  \node[font=\sffamily\tiny, orange!50!black, rotate=90] at (-0.7, 1.75) {Severe};
  \node[font=\sffamily\tiny, red!40!black, rotate=90] at (-0.7, 0.5) {V.~Severe};

  % =====================================================================
  % Ceiling B_max (red dashed, stepping down at each event)
  % =====================================================================
  \draw[red!60!black, thick, densely dashed]
    (0, 7.0) -- (1.5, 7.0)    % pre-event 1
    (1.5, 6.2) -- (3.8, 6.2)  % post-event 1
    (3.8, 5.2) -- (6.0, 5.2)  % post-event 2
    (6.0, 3.9) -- (8.2, 3.9)  % post-event 3 (crosses into moderate)
    (8.2, 2.8) -- (11.0, 2.8);% post-event 4

  % --- Ceiling drop arrows (vertical, at each event) ---
  \foreach \x/\ytop/\ybot in {1.5/7.0/6.2, 3.8/6.2/5.2, 6.0/5.2/3.9, 8.2/3.9/2.8} {
    \draw[->, red!60!black, thick] (\x, \ytop) -- (\x, \ybot);
  }

  % =====================================================================
  % Actual trajectory B(t) (blue, drops sharply then recovers toward ceiling)
  % =====================================================================
  \draw[blue!70!black, very thick]
    % Pre-event 1: at ceiling
    (0, 7.0) -- (1.5, 7.0)
    % Event 1 (infection): sharp drop
    (1.5, 7.0) -- (1.5, 4.5)
    % Recovery toward ceiling 6.2 (doesn't quite reach it)
    .. controls (2.0, 4.8) and (2.8, 5.7) .. (3.8, 6.0)
    % Event 2 (PEM crash): sharp drop
    -- (3.8, 3.5)
    % Recovery toward ceiling 5.2
    .. controls (4.3, 3.8) and (5.2, 4.7) .. (6.0, 5.0)
    % Event 3 (surgery): sharp drop
    -- (6.0, 2.4)
    % Recovery toward ceiling 3.9
    .. controls (6.5, 2.7) and (7.4, 3.4) .. (8.2, 3.7)
    % Event 4 (infection): sharp drop
    -- (8.2, 1.5)
    % Recovery toward ceiling 2.8 (ongoing)
    .. controls (8.8, 1.8) and (9.8, 2.3) .. (11.0, 2.6);

  % =====================================================================
  % Event labels (above, angled)
  % =====================================================================
  \node[font=\sffamily\tiny, orange!80!black, anchor=south, rotate=45]
    at (1.5, 7.3) {Infection};
  \node[font=\sffamily\tiny, orange!80!black, anchor=south, rotate=45]
    at (3.8, 6.5) {PEM crash};
  \node[font=\sffamily\tiny, orange!80!black, anchor=south, rotate=45]
    at (6.0, 5.5) {Surgery};
  \node[font=\sffamily\tiny, orange!80!black, anchor=south, rotate=45]
    at (8.2, 4.2) {Infection};

  % =====================================================================
  % Annotations: delta_k and Delta_k for event 1
  % =====================================================================
  % Irreversible loss (delta_k): gap between old and new ceiling
  \draw[<->, red!60!black, thin] (1.8, 7.0) -- (1.8, 6.2);
  \node[font=\sffamily\tiny, red!60!black, anchor=west] at (1.9, 6.6)
    {$\delta_1$};

  % Total acute drop (Delta_k): from pre-event B to post-event B
  \draw[<->, blue!50!black, thin] (1.2, 7.0) -- (1.2, 4.5);
  \node[font=\sffamily\tiny, blue!50!black, anchor=east] at (1.1, 5.75)
    {$\Delta_1$};

  % --- Legend ---
  \node[draw=gray!40, fill=white, rounded corners=3pt,
        inner sep=4pt, font=\sffamily\scriptsize,
        anchor=north east]
    at (11.0, 7.6) {%
      \begin{tabular}{@{}ll@{}}
        \tikz\draw[blue!70!black, very thick] (0,0) -- (0.5,0);
          & $B(t)$ actual \\[1pt]
        \tikz\draw[red!60!black, thick, densely dashed] (0,0) -- (0.5,0);
          & $B_{\max}$ ceiling \\[1pt]
        \tikz\draw[<->, red!60!black, thin] (0,0) -- (0,0.25);
          & $\delta_k$ (irreversible) \\[1pt]
        \tikz\draw[<->, blue!50!black, thin] (0,0) -- (0,0.25);
          & $\Delta_k$ (total drop) \\
      \end{tabular}%
    };

  % --- Subfigure label ---
  \node[font=\sffamily\bfseries] at (0.6, -0.8) {(a)};

\end{scope}

%% ====================================================================
%% (b) Trajectory Classes (small multiples)
%% ====================================================================
\begin{scope}[shift={(0, -9.5)}]

  % Shared settings for small multiples
  \def\panelw{4.8}
  \def\panelh{3.5}

  % --- Panel (b1): Stable ratchet ---
  \begin{scope}[shift={(0, 0)}]
    \draw[->, thick] (-0.2, 0) -- (\panelw+0.3, 0);
    \draw[->, thick] (0, -0.2) -- (0, \panelh+0.3);
    \node[font=\sffamily\tiny, anchor=north] at (2.4, -0.15) {Time};
    \node[font=\sffamily\tiny, rotate=90, anchor=south] at (-0.35, 1.75) {$B(t)$};

    % Ceiling (equal steps)
    \draw[red!60!black, densely dashed, thin]
      (0, 3.2) -- (1.0, 3.2)
      (1.0, 2.8) -- (2.0, 2.8)
      (2.0, 2.4) -- (3.2, 2.4)
      (3.2, 2.0) -- (4.5, 2.0);

    % Trajectory
    \draw[blue!70!black, thick]
      (0, 3.2) -- (1.0, 3.2) -- (1.0, 2.0)
      .. controls (1.3, 2.2) and (1.7, 2.6) .. (2.0, 2.7)
      -- (2.0, 1.6)
      .. controls (2.3, 1.8) and (2.8, 2.2) .. (3.2, 2.3)
      -- (3.2, 1.2)
      .. controls (3.5, 1.4) and (4.0, 1.8) .. (4.5, 1.9);

    \node[font=\sffamily\scriptsize\bfseries, anchor=north] at (2.4, \panelh+0.2)
      {Stable ($\alpha \approx 0$)};
  \end{scope}

  % --- Panel (b2): Accelerating ratchet ---
  \begin{scope}[shift={(6.2, 0)}]
    \draw[->, thick] (-0.2, 0) -- (\panelw+0.3, 0);
    \draw[->, thick] (0, -0.2) -- (0, \panelh+0.3);
    \node[font=\sffamily\tiny, anchor=north] at (2.4, -0.15) {Time};

    % Ceiling (increasing steps)
    \draw[red!60!black, densely dashed, thin]
      (0, 3.2) -- (1.2, 3.2)
      (1.2, 2.9) -- (2.2, 2.9)
      (2.2, 2.3) -- (3.2, 2.3)
      (3.2, 1.2) -- (4.5, 1.2);

    % Trajectory
    \draw[blue!70!black, thick]
      (0, 3.2) -- (1.2, 3.2) -- (1.2, 2.3)
      .. controls (1.5, 2.5) and (1.9, 2.7) .. (2.2, 2.8)
      -- (2.2, 1.5)
      .. controls (2.5, 1.7) and (2.9, 2.1) .. (3.2, 2.2)
      -- (3.2, 0.4)
      .. controls (3.6, 0.6) and (4.0, 0.9) .. (4.5, 1.0);

    \node[font=\sffamily\scriptsize\bfseries, anchor=north] at (2.4, \panelh+0.2)
      {Accelerating ($\alpha > 0$)};
  \end{scope}

  % Second row labels
  \node[font=\sffamily\bfseries] at (0.6, -0.8) {(b)};

\end{scope}

\begin{scope}[shift={(0, -15.5)}]

  \def\panelw{4.8}
  \def\panelh{3.5}

  % --- Panel (c1): Incomplete-recovery cascade ---
  \begin{scope}[shift={(0, 0)}]
    \draw[->, thick] (-0.2, 0) -- (\panelw+0.3, 0);
    \draw[->, thick] (0, -0.2) -- (0, \panelh+0.3);
    \node[font=\sffamily\tiny, anchor=north] at (2.4, -0.15) {Time};
    \node[font=\sffamily\tiny, rotate=90, anchor=south] at (-0.35, 1.75) {$B(t)$};

    % Ceiling (moderate steps, closely spaced)
    \draw[red!60!black, densely dashed, thin]
      (0, 3.2) -- (0.7, 3.2)
      (0.7, 2.8) -- (1.4, 2.8)
      (1.4, 2.4) -- (2.1, 2.4)
      (2.1, 2.0) -- (2.8, 2.0)
      (2.8, 1.6) -- (3.5, 1.6)
      (3.5, 1.2) -- (4.5, 1.2);

    % Trajectory (barely recovers before next hit)
    \draw[blue!70!black, thick]
      (0, 3.2) -- (0.7, 3.2) -- (0.7, 2.0)
      .. controls (0.85, 2.1) and (1.1, 2.3) .. (1.4, 2.4)
      -- (1.4, 1.3)
      .. controls (1.55, 1.4) and (1.8, 1.6) .. (2.1, 1.7)
      -- (2.1, 0.8)
      .. controls (2.25, 0.9) and (2.5, 1.1) .. (2.8, 1.2)
      -- (2.8, 0.5)
      .. controls (2.95, 0.55) and (3.2, 0.7) .. (3.5, 0.8)
      -- (3.5, 0.3)
      .. controls (3.8, 0.4) and (4.1, 0.6) .. (4.5, 0.7);

    \node[font=\sffamily\scriptsize\bfseries, anchor=north, text width=4cm, align=center]
      at (2.4, \panelh+0.2) {Incomplete Recovery};
  \end{scope}

  % --- Panel (c2): Arrested ratchet ---
  \begin{scope}[shift={(6.2, 0)}]
    \draw[->, thick] (-0.2, 0) -- (\panelw+0.3, 0);
    \draw[->, thick] (0, -0.2) -- (0, \panelh+0.3);
    \node[font=\sffamily\tiny, anchor=north] at (2.4, -0.15) {Time};

    % Two events, then prevention (ceiling stays flat)
    \draw[red!60!black, densely dashed, thin]
      (0, 3.2) -- (1.0, 3.2)
      (1.0, 2.7) -- (2.0, 2.7)
      (2.0, 2.2) -- (4.5, 2.2);

    % Trajectory: two drops, then long recovery toward stable ceiling
    \draw[blue!70!black, thick]
      (0, 3.2) -- (1.0, 3.2) -- (1.0, 1.8)
      .. controls (1.3, 2.0) and (1.7, 2.4) .. (2.0, 2.6)
      -- (2.0, 1.2)
      .. controls (2.5, 1.5) and (3.2, 1.9) .. (4.5, 2.1);

    % Prevention marker
    \draw[green!50!black, thick, densely dotted] (2.0, 0) -- (2.0, 3.5);
    \node[font=\sffamily\tiny, green!50!black, anchor=south, rotate=45]
      at (2.0, 3.5) {Prevention starts};

    \node[font=\sffamily\scriptsize\bfseries, anchor=north, text width=4cm, align=center]
      at (2.4, \panelh+0.2) {Arrested (prevention)};
  \end{scope}

  \node[font=\sffamily\bfseries] at (0.6, -0.8) {(c)};

\end{scope}

\end{tikzpicture}
\caption{The energy ratchet model. (a)~A single patient trajectory showing
baseline functional capacity $B(t)$ (solid blue) and the irreversible ceiling
$B_{\max}$ (dashed red) declining through four damaging events (infections,
PEM crashes, surgery). Each event produces a total acute drop $\Delta_k$
(blue arrow), of which only $\delta_k$ (red arrow) is irreversible.
Recovery between events is partial and ATP-dependent
(Equation~\ref{eq:ratchet-interevent}). Horizontal bands indicate clinical
severity levels (Equation~\ref{eq:ratchet-severity}).
(b)~Two of four trajectory classes: \emph{stable} ratchet with equal ceiling
steps ($\alpha \approx 0$) and \emph{accelerating} ratchet with damage
sensitisation ($\alpha > 0$).
(c)~\emph{Incomplete-recovery cascade} (frequent events prevent recovery)
and \emph{arrested ratchet} (event prevention preserves ceiling).
See Equations~\ref{eq:ratchet-jump-ceiling}--\ref{eq:ratchet-irreversible-loss}
for the formal jump map.}
\label{fig:energy-ratchet}
\end{figure}
